\subsection{平稳随机过程}
\subsubsection{平稳性和遍历性}
\textbf{狭义平稳 (严平稳)}\quad 统计特性与时间起点无关.
\begin{itemize}
    \item $f_X(x;t)$ 与 $t$ 无关;
    \item $f_X(x_1,x_2;t,t+\tau)$ 仅与 $\tau$ 有关.
\end{itemize}

\textbf{广义平稳}
\begin{itemize}
    \item $a(t)$ 与 $t$ 无关;
    \item $R(t,t+\tau)$ 仅与 $\tau$ 有关.
\end{itemize}

\textbf{各态历经性 (遍历性)}\quad (对于\textit{平稳过程}而言) 任何一个样本函数 $x(t)$ 都经历了所有可能状态, 即其相关指标的时间均值与统计均值相等.
\begin{itemize}
    \item $\overline{a(t)}=\displaystyle\lim_{T\rightarrow\infty}\frac{1}{T}\int_{-T/2}^{T/2}x(t)\mathrm{d}t=a$;
    \item $\overline{R(\tau)}=\displaystyle\lim_{T\rightarrow\infty}\frac{1}{T}\int_{-T/2}^{T/2}x(t)x(t+\tau)\mathrm{d}t=R(\tau)$.
\end{itemize}

\begin{exampleprob}[各态历经性的判断]
    设相位随机的正弦随机过程为 $X(t)=A\cos(\omega_ct+\theta)$, 其中 $A$ 和 $\omega_c$ 为常数, $\theta$ 是在 $(0,2\pi)$ 内均匀分布的随机变量. 判断 $X(t)$ 是否具有各态历经性.

    \begin{solution}
        因为 $X(t)=A(\cos\omega_ct\cos\theta-\sin\omega_ct\sin\theta)$, 所以
        \begin{equation} \label{eq:3.2 eprob 1}
            \begin{aligned}
                a(t) & =E[A(\cos\omega_ct\cos\theta-\sin\omega_ct\sin\theta)]   \\
                     & =A\cos\omega_ctE[\cos\theta]-A\sin\omega_ctE[\sin\theta] \\
                     & =0,
            \end{aligned}
        \end{equation}
        其中利用了
        \begin{equation*}
            E[\cos\theta]=\int_{-\infty}^{+\infty}\cos\theta\cdot\frac{1}{2\pi}\mathrm{d}\theta=0.
        \end{equation*}
        同理可得 $E[\sin\theta]=0$.

        因为
        \begin{equation*}
            \begin{aligned}
                X(t)X(t+\tau) & =A\cos(\omega_ct+\theta)\cdot A\cos(\omega_c(t+\tau)+\theta)                                            \\
                              & =\frac{A^2}{2}[\cos(\omega_c(2t+\tau)+2\theta)+\cos\omega_c\tau]                                        \\
                              & =\frac{A^2}{2}[\cos(\omega_c(2t+\tau))\cos2\theta-\sin(\omega_c(2t+\tau))\sin2\theta+\cos\omega_c\tau],
            \end{aligned}
        \end{equation*}
        且容易得到 $E[\cos2\theta]=E[\sin2\theta]=0$, 所以
        \begin{equation} \label{eq:3.2 eprob 2}
            R_X(t,t+\tau)=E[X(t)X(t+\tau)]=\frac{A^2}{2}\cos\omega_c\tau.
        \end{equation}

        根据式 (\ref{eq:3.2 eprob 1}) 和 (\ref{eq:3.2 eprob 2}), $X(t)$ 是平稳过程.

        因为
        \begin{equation*}
            \begin{aligned}
                  & \int_{-T/2}^{T/2}A\cos(\omega_ct+\theta)\mathrm{d}t                                                                                                                    \\
                = & \left.\frac{A}{\omega_c}\sin(\omega_ct+\theta)\right|_{-T/2}^{T/2}                                                                                                     \\
                = & \frac{A}{\omega_c}\left[\sin\left(\frac{\omega_cT}{2}+\theta\right)-\sin\left(\frac{\omega_cT}{2}+\theta\right)\right]                                                 \\
                = & \frac{A}{\omega_c}\left(\sin\frac{\omega_cT}{2}\cos\theta+\cos\frac{\omega_cT}{2}\sin\theta+\sin\frac{\omega_cT}{2}\cos\theta-\cos\frac{\omega_cT}{2}\sin\theta\right) \\
                = & \frac{2A}{\omega_c}\sin\frac{\omega_cT}{2}\cos\theta,
            \end{aligned}
        \end{equation*}
        所以
        \begin{equation} \label{eq:3.2 eprob 3}
            \begin{aligned}
                \overline{a(t)} & =\lim_{T\rightarrow\infty}\frac{1}{T}\int_{-T/2}^{T/2}X(t)\mathrm{d}t           \\
                                & =\lim_{T\rightarrow\infty}\frac{2A}{\omega_cT}\sin\frac{\omega_cT}{2}\cos\theta \\
                                & =0.
            \end{aligned}
        \end{equation}

        因为
        \begin{equation*}
            \begin{aligned}
                \overline{R(t,t+\tau)} & =\lim_{T\rightarrow\infty}\frac{1}{T}\int_{-T/2}^{T/2}A\cos(\omega_ct+\theta)\cdot A\cos(\omega_c(t+\tau)+\theta)\mathrm{d}t            \\
                                       & =\lim_{T\rightarrow\infty}\frac{A^2}{2T}\int_{-T/2}^{T/2}[\cos\omega_c\tau\mathrm{d}t+\cos(2\omega_ct+\omega_c\tau+2\theta)]\mathrm{d}t \\
                                       & =\frac{A^2}{2}\cos\omega_c\tau,
            \end{aligned}
        \end{equation*}
        以及式 (\ref{eq:3.2 eprob 3}), 所以 $X(t)$ 是各态历经过程.
    \end{solution}
\end{exampleprob}

\subsubsection{自相关函数的性质}
\begin{itemize}
    \item 偶函数;
    \item \underline{平均功率}\quad $S=E[X^2(t)]=R(0)$;
    \item \underline{直流功率}\quad $a^2=E^2[X(t)]=R(\infty)$;
    \item \underline{交流功率}\quad $\sigma^2=R(0)-R(\infty)$;
    \item $R(0)\geq|R(\tau)|$.
\end{itemize}

其中, 最后一项性质, 可由 $E[(X(t)\pm X(t+\tau))^2]\geq 0$ 左边展开得到 $R(0)\pm R(\tau)\geq 0$.

\subsubsection{功率谱密度}
对于确定功率信号 $f(t)$, \textbf{功率谱密度}定义为
\begin{equation}
    P_f(\omega)=\lim_{T\rightarrow\infty}\frac{|F_T(\omega)|^2}{T},
\end{equation}
其中 $F_T(\omega)$ 表示 $f(t)$ 的截短函数 $f_T(t)$ 对应的频谱函数.

不同样本函数具有不同的功率谱密度. 一个随机过程的功率谱密度定义为所有样本函数的功率谱的\textit{统计平均}:
\begin{equation}
    P_X(\omega)=E[P_f(\omega)]=\lim_{T\rightarrow\infty}\frac{E(|F_T(\omega)|^2)}{T}.
\end{equation}

\textbf{维纳-辛钦定理}
\begin{equation} \label{eq:3.2 wiener}
    P_X(\omega)=\mathcal{F}[R(\tau)].
\end{equation}

\textbf{平稳过程的功率谱性质}
\begin{itemize}
    \item $R(0)=\displaystyle\int_{-\infty}^{\infty}P_X(\omega)\mathrm{d}\omega$;
    \item 各态历经过程的任一样本函数的功率谱即为随机过程的功率谱;
    \item $P_X(\omega)\geq 0$, $P_X(\omega)=P_X(-\omega)$.
\end{itemize}
